\section{Discussion}
\label{sec:discussion}

The results show two exposures producing responses of different kind. Both act on the same object, the cost of traversing an edge. A continuous exposure applies a finite factor, above one to penalise the traverse or below one to mitigate it, as heat does, and the effect shows up as a shift in the distribution of travel times. A discrete exposure applies an infinite factor, whose practical effect is the removal of the edge, as landsliding does, and the effect shows up as a fragmented network. A composite exposure inherits both without averaging them. That separation is what the two-stage design was meant to deliver, and it holds in Curitiba and Porto Alegre despite their different natures, one with a single thermal layer and the other with two vector layers applied in sequence. The magnitudes differ as much as the signatures. Heat moves the city average from 12.6 to 13.1 minutes, a median cell gain of half a minute, while flooding moves it from 15.1 to 20.7 with a 90th percentile of 46 minutes and landsliding takes 619 cells out of the reachable set entirely. Both exposures are modest against the three to four fold spread between core and periphery that the typical scenario already contains, which places them where a single average cannot separate them.

The socioeconomic reading is less direct than the aggregate numbers suggest. The SDI is used here as a stratifier of the population, not as an explanatory variable of the mechanism, since neither the network nor the penalties read census data. Any association between the SDI and the change in access comes from the geography of the city, from where the exposed segments sit relative to where the population with fewer resources lives. In Porto Alegre the association is consistently positive and sanitation carries most of it, which is what a compound exposure produces in a city whose service provision already follows the terrain. In Curitiba the association weakens once vegetation enters, and the component that remains moves to income. Where the surface was hot, sanitation was also poor, so the sanitation axis partly measured the heat itself rather than its cause. Removing part of the heat also removes the part of the correlation that the heat carried, and sanitation falls from 0.350 to 0.098.

That design choice has a cost. A composite of three equally weighted axes is a summary of deprivation, not a measurement of exposure, and the regression shows that the axes are not interchangeable across cities and scenarios. Porto Alegre concentrates the relation in one axis under \texttt{flood} and collapses all three under \texttt{landslide}. Curitiba redistributes it when vegetation changes the surface. Two cells can share the same SDI while one owes it to sanitation and the other to income, so the index cannot tell which deficit drives the loss. Any agreement it shows with the change in access is an agreement with whichever axis happens to dominate in that cell, not with deprivation as a whole. The SDI therefore belongs to the population side of the analysis and not to the explanation of the mechanism, and the regression on its three axes, rather than on the composite, is where the mechanism enters.

Several limits bound the magnitudes. Surface temperature is a proxy for air temperature, and the gap between them is spatially systematic, larger where the urban heat island is stronger, which is exactly where the penalty matters. The conversion from temperature to walking speed comes from a laboratory study of young healthy adults, so the coefficient describes a sensitivity and not a measured population. The vegetation factor comes from a study of perceived travel time under street trees, applied here as a uniform reduction wherever tree cover is present, without conditioning on heat. For Porto Alegre the flood penalty takes its high-class factor from a laboratory measurement and sets the two lower classes by a methodological choice on the same scale, without direct empirical support for those levels. Landsliding is modelled as a blockage on the assumption that the event has occurred. The network itself comes from OpenStreetMap, so the results describe the mapped pedestrian infrastructure and its absence of sidewalks is not represented.

The indicators measure potential access, not realised travel. A model of shortest paths assumes that every pedestrian knows the network and chooses the least penalised route, and that displacement continues on foot under conditions in which a real traveller might cancel the trip, change the destination or use another mode. Interdiction and penalisation are static, applied to a single event, while in a real flood the water rises and recedes and in a real landslide the alternative routes are themselves rebuilt. What the scenarios answer is where a given exposure removes or degrades access within an existing network, and which population sits at the end of the affected routes. Whether the magnitudes correspond to an observed event is a question that requires calibration against field data, which the demonstrative scenarios do not attempt.

Two recommendations for measurement follow from this reading rather than from a test performed here. The first is that the city average is the wrong scale for the decision. A municipality can lose little access on average and still concentrate the loss on the tracts with the least capacity to absorb it, as the local indicators of Section~\ref{subsec:results-lisa} show, and the aggregate hides that concentration. The second is that how an edge is affected carries as much information as how many edges are. The same input layer, read as a penalisation and as an interdiction, produces two different spatial signatures, and the framework developed here treats them as two settings of one parameter rather than as two separate models, so the comparison costs a change of factor and not a change of pipeline.
